\documentclass[11pt,reqno]{amsart}

% ==================================================
% 1. Packages
% ==================================================

\usepackage{amsmath,amssymb,amsthm,mathtools}
\usepackage{mathrsfs}
\usepackage{graphicx}
\usepackage{enumitem}
\usepackage{tikz}
\usepackage[hidelinks]{hyperref}
\usepackage[nameinlink,noabbrev]{cleveref}
\usepackage{mathrsfs} 
% ==================================================
% 2. Mathematical operators
% ==================================================

\DeclareMathOperator{\Aut}{Aut}
\DeclareMathOperator{\Stab}{Stab}
\DeclareMathOperator{\Fix}{Fix}
\DeclareMathOperator{\rank}{rank}
\DeclareMathOperator{\im}{im}
\DeclareMathOperator{\Ker}{Ker}
\DeclareMathOperator{\PLC}{PLC}

\DeclareMathOperator{\Inv}{Inv}
% ==================================================
% 3. Mathematical symbols
% ==================================================

\newcommand{\R}{\mathbb{R}}
\newcommand{\C}{\mathbb{C}}
\newcommand{\Z}{\mathbb{Z}}
\newcommand{\N}{\mathbb{N}}

\newcommand{\norm}[1]{\left\lVert #1 \right\rVert}
\newcommand{\abs}[1]{\left\lvert #1 \right\rvert}
\newcommand{\pairing}[2]{\left\langle #1,#2 \right\rangle}

% Coxeter group notation (optional)

\newcommand{\simpleRoots}{\Pi}
\newcommand{\limitRoots}{\mathcal{E}}
\newcommand{\imaginaryCone}{\mathcal{I}}
\newcommand{\TitsCone}{\mathcal{T}}
\newcommand{\depth}{\operatorname{dp}}

% ==================================================
% 4. Theorem environments
% ==================================================

\theoremstyle{plain}

\newtheorem{theorem}{Theorem}[section]
\newtheorem{lemma}[theorem]{Lemma}
\newtheorem{proposition}[theorem]{Proposition}
\newtheorem{corollary}[theorem]{Corollary}

\theoremstyle{definition}

\newtheorem{definition}[theorem]{Definition}
\newtheorem{example}[theorem]{Example}
\newtheorem{notation}[theorem]{Notation}

\theoremstyle{remark}

\newtheorem{remark}[theorem]{Remark}

% ==================================================
% 5. Equation numbering
% ==================================================

\numberwithin{equation}{section}

% ==================================================
% 6. Hyper-reference names
% ==================================================

\crefname{theorem}{Theorem}{Theorems}
\crefname{lemma}{Lemma}{Lemmas}
\crefname{proposition}{Proposition}{Propositions}
\crefname{corollary}{Corollary}{Corollaries}
\crefname{definition}{Definition}{Definitions}
\crefname{equation}{Equation}{Equations}
\crefname{section}{Section}{Sections}
\DeclareMathOperator{\tr}{tr}
\DeclareMathOperator{\supp}{supp}
\DeclareMathOperator{\Pos}{Pos}
\DeclareMathOperator{\rad}{rad}
\DeclareMathOperator{\dom}{dom}

% ==================================================
% 7. Title and author
% ==================================================

\title[Counterexamples]
{Counterexamples}

\author{Haoquan Wang}

\address{
School of Mathematical Science,
Peking University,
Peking, China
}

\email{2300010755@stu.pku.edu.cn}

% Optional: AMS classification and keywords


\keywords{
Coxeter groups, Fushcian groups, limit roots,
dominance order
}

% ==================================================
% 8. Document
% ==================================================

\begin{document}

% Abstract

\begin{abstract}

\end{abstract}

\maketitle
\section{Counterexample}
\begin{proposition}
Let $(W,S)$ be the Coxeter system with simple roots
$\Pi=\{\alpha_1,\ldots,\alpha_6\}$ and Gram matrix
\[
G=
\begin{pmatrix}
1&-3&0&0&0&0\\
-3&1&-3&0&0&0\\
0&-3&1&-3&0&0\\
0&0&-3&1&-3&0\\
0&0&0&-3&1&-3\\
0&0&0&0&-3&1
\end{pmatrix}.
\]
Then $B$ has signature $(3,3)$, and there exists an infinite strict
dominance chain $(\beta_n)_{n\geq1}$ in $\Phi^+$ such that
$(\widehat{\beta_n})_{n\geq1}$ has exactly two distinct accumulation
points.
\end{proposition}

\begin{proof}
We use the convention that $\beta\succ\alpha$ means that $\beta$
strictly dominates $\alpha$. Write
\[
h\left(\sum_{i=1}^6 c_i\alpha_i\right)=\sum_{i=1}^6 c_i,
\qquad
\widehat{z}=\frac{z}{h(z)},
\]
whenever $h(z)>0$. All vector and operator norms below are the
$\ell^1$ norms relative to $\Pi$.

The eigenvalues of $G$ are
\[
1-6\cos\frac{k\pi}{7},
\qquad 1\leq k\leq6.
\]
Exactly three are positive and three are negative. Thus $B$ has
signature $(3,3)$.

Put
\[
g=s_1s_2,\qquad t=s_5s_6,\qquad
r=3+2\sqrt2,\qquad
\lambda=r^2=17+12\sqrt2,
\]
and define
\[
x=\frac{r\alpha_1+\alpha_2}{r+1},
\qquad
y=\frac{r\alpha_5+\alpha_6}{r+1}.
\]
These are limit roots of the two standard infinite dihedral
subgroups. In particular,
\[
x,y\in E,\qquad B(x,y)=0,\qquad \|x-y\|_1=2.
\]

For $m\geq0$, let
\[
u_m=g^m\alpha_3,\qquad v_m=t^m\alpha_4.
\]
The root $u_m$ is supported on $\{1,2,3\}$, with
$\alpha_3$-coefficient equal to $1$. Similarly, $v_m$ is supported
on $\{4,5,6\}$, with $\alpha_4$-coefficient equal to $1$.
The positive--negative dichotomy for real roots therefore implies
that $u_m,v_m\in\Phi^+$.

The only nonzero pairing between these two supports is
$B(\alpha_3,\alpha_4)=-3$. Consequently,
\[
B(u_m,v_\ell)=-3,\qquad
B(u_m,y)=0,\qquad B(v_\ell,x)=0
\qquad(m,\ell\geq0).
\tag{1}
\]
Moreover, $gx=\lambda x$ and $ty=\lambda y$, so
\[
B(u_m,x)=-\frac{3\lambda^{-m}}{r+1}<0,
\qquad
B(v_m,y)=-\frac{3r\lambda^{-m}}{r+1}<0.
\tag{2}
\]

We first record quantitative convergence estimates. Direct
diagonalization gives
\[
u_m=\frac{3}{16}\lambda^m(r\alpha_1+\alpha_2)+R_m,
\]
where
\[
R_m=
\alpha_3-\frac98\alpha_1-\frac38\alpha_2
+\frac{3}{16}\lambda^{-m}(r^{-1}\alpha_1+\alpha_2),
\]
and
\[
v_m=\frac{3}{16r}\lambda^m(r\alpha_5+\alpha_6)+T_m,
\]
where
\[
T_m=
\alpha_4-\frac38\alpha_5-\frac98\alpha_6
+\frac{3r}{16}\lambda^{-m}(r^{-1}\alpha_5+\alpha_6).
\]
Since $1<r<6$, these remainders satisfy
\[
\|R_m\|_1\leq3,\qquad \|T_m\|_1\leq4.
\]
For nonzero vectors $a,b$ in the positive cone,
\[
\left\|\frac{a}{h(a)}-\frac{b}{h(b)}\right\|_1
\leq\frac{2\|a-b\|_1}{h(b)}.
\]
Applying this inequality to the preceding decompositions yields
\[
\|\widehat{u_m}-x\|_1\leq50\lambda^{-m},
\qquad
\|\widehat{v_m}-y\|_1\leq50\lambda^{-m}.
\tag{3}
\]

Now set
\[
M_n=10^n,\qquad
\gamma_n=
\begin{cases}
u_{M_n},& n\text{ odd},\\
v_{M_n},& n\text{ even},
\end{cases}
\]
and define
\[
w_0=1,\qquad
w_n=s_{\gamma_1}\cdots s_{\gamma_n},
\qquad
\beta_n=w_{n-1}\gamma_n.
\tag{4}
\]

We verify the dominance relations before studying convergence.
If $a,b\in\Phi^+$ satisfy $B(a,b)=-3$, then
\[
s_ab=b+6a\in\Phi^+,
\qquad B(s_ab,a)=3.
\]
By the standard dominance comparability criterion, $s_ab$ and $a$
are comparable. The relation $a\succ s_ab$ is impossible:
applying $s_a$ sends $a$ to the negative root $-a$ but sends
$s_ab$ to the positive root $b$. Hence
\[
s_ab\succ a.
\tag{5}
\]
By (1), $B(\gamma_n,\gamma_{n+1})=-3$. Applying $w_{n-1}$ to (5)
and using $W$-invariance of dominance gives
\[
\beta_{n+1}
=w_{n-1}s_{\gamma_n}\gamma_{n+1}
\succ w_{n-1}\gamma_n
=\beta_n.
\]
Since $\beta_1\in\Phi^+$, every $\beta_n$ is positive. The chain is
strict, and
\[
B(\beta_n,\beta_{n+1})=3.
\tag{6}
\]

Define the moving limit roots
\[
X_n=\widehat{w_nx},\qquad
Y_n=\widehat{w_ny}.
\]
These are well defined because $E\subseteq U^*$ and hence
$wx,wy$ belong to the positive cone for every $w\in W$.

We claim that
\[
\|\widehat{\beta_n}-X_{n-1}\|_1<2^{-n-4}
\quad\text{if }n\text{ is odd},
\tag{7}
\]
and
\[
\|\widehat{\beta_n}-Y_{n-1}\|_1<2^{-n-4}
\quad\text{if }n\text{ is even}.
\tag{8}
\]
To prove this, observe that if $p,q$ are normalized positive vectors
and $wp,wq$ are positive, then
\[
\|\widehat{wp}-\widehat{wq}\|_1
\leq
2\|w\|_1\|w^{-1}\|_1\|p-q\|_1.
\tag{9}
\]
Indeed, the normalization inequality above gives a bound with
denominator $h(wq)$, and
\[
h(wq)=\|wq\|_1\geq\|w^{-1}\|_1^{-1}.
\]

The matrices of $g$ and $t$ satisfy
\[
\|g\|_1,\|t\|_1\leq43.
\]
Consequently, $\|\gamma_j\|_1\leq43^{M_j}$. The reflection formula
and the bound $\max_{i,j}|G_{ij}|\leq3$ imply
\[
\|s_{\gamma_j}\|_1
\leq1+6\|\gamma_j\|_1^2
\leq7\cdot43^{2M_j}.
\]
Since each reflection is its own inverse,
\[
\|w_{n-1}\|_1\|w_{n-1}^{-1}\|_1
\leq49^{n-1}43^{4\sum_{j<n}M_j}.
\tag{10}
\]
Combining (3), (9), and (10), the left-hand side of either
(7) or (8) is at most
\[
100\cdot49^{n-1}
43^{4\sum_{j<n}M_j}\lambda^{-M_n}.
\]
Now
\[
\sum_{j<n}M_j=\frac{10^n-10}{9},
\qquad
43<\lambda^2,\qquad49<\lambda^2.
\]
Therefore this bound is at most
\[
100\lambda^{\,2(n-1)-(10^n+80)/9}
\leq100\lambda^{-n-9}
<2^{-n-4}.
\]
Here the middle inequality follows from
$10^n\geq27n-17$ for $n\geq1$, and the last follows from
$\lambda>33$. This proves (7) and (8).

Suppose first that $n$ is odd. By (1),
$s_{\gamma_n}y=y$, and therefore
\[
Y_n=Y_{n-1}.
\]
By (2), the scalar $c_n=-2B(\gamma_n,x)$ is positive, and
\[
w_nx=w_{n-1}x+c_n\beta_n.
\]
All vectors on the right are positive. Thus $X_n$ is a convex
combination of $X_{n-1}$ and $\widehat{\beta_n}$. By (7),
\[
\|X_n-X_{n-1}\|_1<2^{-n-4}.
\tag{11}
\]
Similarly, for even $n$,
\[
X_n=X_{n-1},
\qquad
\|Y_n-Y_{n-1}\|_1<2^{-n-4}.
\tag{12}
\]

Since
\[
\sum_{n=1}^{\infty}2^{-n-4}=\frac1{16},
\]
both $(X_n)$ and $(Y_n)$ are Cauchy sequences. Write
\[
X_n\longrightarrow X,\qquad Y_n\longrightarrow Y.
\]
Their limits satisfy
\[
\|X-x\|_1\leq\frac1{16},
\qquad
\|Y-y\|_1\leq\frac1{16}.
\]
It follows that
\[
\|X-Y\|_1
\geq\|x-y\|_1-\|X-x\|_1-\|Y-y\|_1
\geq\frac{15}{8}>0.
\]
In particular, $X\neq Y$.

Finally, (7) and (8) imply
\[
\widehat{\beta_{2k-1}}\longrightarrow X,
\qquad
\widehat{\beta_{2k}}\longrightarrow Y.
\]
Thus the normalized strict dominance chain has exactly the two
accumulation points $X$ and $Y$. Both belong to $E$. Moreover,
\[
B(X_n,Y_n)
=\frac{B(x,y)}{h(w_nx)h(w_ny)}=0,
\]
so continuity gives $B(X,Y)=0$.
\end{proof}

\begin{proposition}\label{prop:counterexample-convE-Q}
There exists an irreducible based root system of rank five and
signature $(3,2)$ for which
\[
E \subsetneq \operatorname{conv}(E)\cap Q.
\]
More precisely, let
\[
V=\bigoplus_{i=1}^{5}\mathbb{R}\alpha_i
\]
and let $B$ be the symmetric bilinear form whose Gram matrix,
with respect to the basis $(\alpha_1,\ldots,\alpha_5)$, is
\[
G=
\begin{pmatrix}
1&-2&0&0&-2\\
-2&1&0&0&-2\\
0&0&1&-2&-2\\
0&0&-2&1&-2\\
-2&-2&-2&-2&1
\end{pmatrix}.
\]
Let $W=\langle s_1,\ldots,s_5\rangle$, where
\[
s_i(v)=v-2B(v,\alpha_i)\alpha_i,
\]
and let
\[
\Phi=W\{\alpha_1,\ldots,\alpha_5\},
\qquad
P=\sum_{i=1}^{5}\mathbb{R}_{\geq 0}\alpha_i,
\qquad
\Phi^+=\Phi\cap P.
\]
Define
\[
h\left(\sum_{i=1}^{5}c_i\alpha_i\right)=\sum_{i=1}^{5}c_i,
\qquad
\widehat{v}=\frac{v}{h(v)},
\qquad
E=\operatorname{Acc}\{\widehat{\rho}:\rho\in\Phi^+\},
\]
and put
\[
Q=\{v\in V:B(v,v)=0\}.
\]
Set
\[
\tau=2-\sqrt{3},
\qquad
x=\frac{\alpha_1+\tau\alpha_2}{1+\tau},
\qquad
y=\frac{\alpha_3+\tau\alpha_4}{1+\tau}.
\]
Then
\[
x,y\in E,
\qquad
B(x,y)=0,
\qquad
\frac{x+2y}{3}\notin E.
\]
\end{proposition}

\begin{proof}
The off-diagonal entries of $G$ belong to $\{0,-2\}$, so $G$
defines a based root system: an entry $0$ corresponds to
$m_{ij}=2$, whereas an entry $-2$ corresponds to $m_{ij}=\infty$.
Its Coxeter graph is connected, since the fifth vertex is
joined to each of the other four vertices. Hence the root
system is irreducible.

Write
\[
x_+=x,
\qquad
x_-=\frac{\tau\alpha_1+\alpha_2}{1+\tau},
\]
and
\[
y_+=y,
\qquad
y_-=\frac{\tau\alpha_3+\alpha_4}{1+\tau}.
\]
These are the limit roots of the two non-affine dihedral
root subsystems supported on
$\{\alpha_1,\alpha_2\}$ and $\{\alpha_3,\alpha_4\}$.

For completeness, the matrix of $s_1s_2$ on
$\operatorname{span}\{\alpha_1,\alpha_2\}$ is
\[
\begin{pmatrix}
15&-4\\
4&-1
\end{pmatrix}.
\]
Its eigenvectors $x_+$ and $x_-$ have respective eigenvalues
\[
\tau^{-2}=7+4\sqrt{3},
\qquad
\tau^2=7-4\sqrt{3}.
\]
Consequently,
\[
\widehat{(s_1s_2)^n\alpha_1}\longrightarrow x_+.
\]
The corresponding argument for $s_3s_4$ gives
\[
\widehat{(s_3s_4)^n\alpha_3}\longrightarrow y_+.
\]
Thus $x_+,y_+\in E$.

Direct calculation gives
\[
B(x_+,x_+)=B(x_-,x_-)
=B(y_+,y_+)=B(y_-,y_-)=0,
\]
\[
B(x_+,x_-)=B(y_+,y_-)=-1,
\]
and
\[
B(x_\varepsilon,y_\delta)=0
\qquad
(\varepsilon,\delta\in\{+,-\}).
\]
Let
\[
\eta=\alpha_5-2\alpha_1-2\alpha_2-2\alpha_3-2\alpha_4.
\]
Then
\[
\eta\perp_B
\operatorname{span}\{x_+,x_-,y_+,y_-\},
\qquad
B(\eta,\eta)=17.
\]
The vectors $(x_+,x_-,y_+,y_-,\eta)$ form a basis of $V$,
and the Gram matrix in this basis is
\[
\begin{pmatrix}
0&-1&0&0&0\\
-1&0&0&0&0\\
0&0&0&-1&0\\
0&0&-1&0&0\\
0&0&0&0&17
\end{pmatrix}.
\]
Therefore $B$ has signature $(3,2)$.

For $v=\sum_{i=1}^{5}c_i\alpha_i$, define linear forms
\[
\begin{aligned}
A_+(v)
&=\frac{1+\sqrt{3}}{2}c_1
  -\frac{\sqrt{3}-1}{2}c_2+2c_5,\\
A_-(v)
&=-\frac{\sqrt{3}-1}{2}c_1
  +\frac{1+\sqrt{3}}{2}c_2+2c_5,\\
C_+(v)
&=\frac{1+\sqrt{3}}{2}c_3
  -\frac{\sqrt{3}-1}{2}c_4+2c_5,\\
C_-(v)
&=-\frac{\sqrt{3}-1}{2}c_3
  +\frac{1+\sqrt{3}}{2}c_4+2c_5.
\end{aligned}
\]
These are the first four coordinates in the preceding basis:
\[
v=
A_+(v)x_++A_-(v)x_-
+C_+(v)y_++C_-(v)y_-+c_5\eta.
\]

We claim that every positive root $\rho$ satisfies
\begin{equation}\label{eq:root-gap-inequality}
\bigl(20A_+(\rho)-9C_+(\rho)\bigr)
\bigl(3A_+(\rho)-2C_+(\rho)\bigr)\geq 0.
\end{equation}
We first establish a useful description of positive roots.

Since the simple reflections have integer matrices in the
simple-root basis, every root has integer coordinates.
Let $\rho=\sum_i c_i\alpha_i$ be a positive root that is
not simple. Since
\[
1=B(\rho,\rho)=\sum_i c_i B(\rho,\alpha_i),
\]
there exists $i$ with $B(\rho,\alpha_i)>0$.
A simple reflection preserves all positive roots other than
its corresponding simple root. Hence
\[
s_i\rho\in\Phi^+,
\qquad
h(s_i\rho)=h(\rho)-2B(\rho,\alpha_i)<h(\rho).
\]
The height is a positive integer, so repeating this procedure
eventually reaches a simple root.

Reversing the resulting sequence shows that every positive
root can be obtained from a simple root by a sequence of
simple reflections in which each step strictly increases
one simple-root coordinate and leaves the others unchanged.

Suppose first that the fifth coordinate of $\rho$ is positive.
Choose such a coordinate-increasing sequence.
If the sequence contains $s_5$, let $v$ be the root immediately
after its last occurrence.
If it contains no $s_5$, then it must start at $\alpha_5$;
in this case put $v=\alpha_5$.
In either case,
\[
\rho=uv
\qquad\text{for some}\qquad
u\in W_{\{1,2,3,4\}}.
\]

Write
\[
v=\sum_{i=1}^{4}c_i\alpha_i+d\alpha_5,
\qquad
S=c_1+c_2+c_3+c_4.
\]
We have
\begin{equation}\label{eq:seed-coordinate-bound}
d\geq 2S.
\end{equation}
Indeed, if $v=\alpha_5$, this is immediate.
Otherwise, let $d_0$ be the fifth coordinate immediately
before the last occurrence of $s_5$.
Since that step increases the fifth coordinate,
\[
d_0-2S
=
B\left(\sum_{i=1}^{4}c_i\alpha_i+d_0\alpha_5,\alpha_5\right)
<0.
\]
The fifth coordinate after reflection is
\[
d=4S-d_0>2S,
\]
which proves \eqref{eq:seed-coordinate-bound}.

Since $c_i\geq 0$ and $S\leq d/2$, the definitions of the
four linear forms imply
\[
\frac{9-\sqrt{3}}{4}\,d
\leq
A_\varepsilon(v),C_\varepsilon(v)
\leq
\frac{9+\sqrt{3}}{4}\,d
\qquad
(\varepsilon\in\{+,-\}).
\]
In particular, these four values are strictly positive.
Put
\[
R=\frac{9+\sqrt{3}}{9-\sqrt{3}}.
\]
Then
\begin{equation}\label{eq:seed-ratio-bound}
R^{-1}
\leq
\frac{A_\varepsilon(v)}{C_\delta(v)}
\leq R
\qquad
(\varepsilon,\delta\in\{+,-\}).
\end{equation}
Moreover,
\begin{equation}\label{eq:elementary-gap-bounds}
R<\frac32,
\qquad
0<\tau<\frac3{10}.
\end{equation}

The action of the first two simple reflections on the
coordinates $(A_+,A_-)$ is
\[
\begin{pmatrix}
A_+(s_1v)\\
A_-(s_1v)
\end{pmatrix}
=
\begin{pmatrix}
0&\tau^{-1}\\
\tau&0
\end{pmatrix}
\begin{pmatrix}
A_+(v)\\
A_-(v)
\end{pmatrix},
\]
and
\[
\begin{pmatrix}
A_+(s_2v)\\
A_-(s_2v)
\end{pmatrix}
=
\begin{pmatrix}
0&\tau\\
\tau^{-1}&0
\end{pmatrix}
\begin{pmatrix}
A_+(v)\\
A_-(v)
\end{pmatrix}.
\]
They fix the coordinates $(C_+,C_-)$ and the vector $\eta$.
The reflections $s_3,s_4$ act by the same two matrices on
$(C_+,C_-)$ and fix $(A_+,A_-)$ and $\eta$.

It follows that there exist integers $k,\ell$ and signs
$\varepsilon,\delta\in\{+,-\}$ such that
\[
A_+(\rho)=\tau^k A_\varepsilon(v),
\qquad
C_+(\rho)=\tau^\ell C_\delta(v).
\]
In particular, $A_+(\rho),C_+(\rho)>0$, and
\[
\frac{A_+(\rho)}{C_+(\rho)}
\in
\tau^{k-\ell}[R^{-1},R].
\]
If $k-\ell\leq 0$, then
\[
\frac{A_+(\rho)}{C_+(\rho)}
\geq R^{-1}>\frac23.
\]
If $k-\ell\geq 1$, then
\[
\frac{A_+(\rho)}{C_+(\rho)}
\leq \tau R<\frac9{20}.
\]
Thus the ratio cannot lie between $9/20$ and $2/3$.
This proves \eqref{eq:root-gap-inequality} whenever the
fifth coordinate of $\rho$ is positive.

Suppose now that the fifth coordinate of $\rho$ is zero.
Along a coordinate-increasing sequence producing $\rho$,
the fifth coordinate must remain zero. Hence the sequence
does not use $s_5$ and starts at one of
$\alpha_1,\ldots,\alpha_4$.
Because the two blocks $\{1,2\}$ and $\{3,4\}$ are orthogonal,
the resulting root is supported entirely on one of them.

If $\rho$ is supported on $\{\alpha_1,\alpha_2\}$, then
$C_+(\rho)=0$, and the left-hand side of
\eqref{eq:root-gap-inequality} is
\[
60A_+(\rho)^2\geq 0.
\]
If $\rho$ is supported on $\{\alpha_3,\alpha_4\}$, then
$A_+(\rho)=0$, and the same expression equals
\[
18C_+(\rho)^2\geq 0.
\]
This completes the proof of
\eqref{eq:root-gap-inequality} for every positive root.

Define the continuous homogeneous quadratic function
\[
F(v)=
\bigl(20A_+(v)-9C_+(v)\bigr)
\bigl(3A_+(v)-2C_+(v)\bigr).
\]
For every positive root $\rho$,
\[
F(\widehat{\rho})=\frac{F(\rho)}{h(\rho)^2}\geq 0.
\]
Passing to accumulation points gives
\[
F(e)\geq 0
\qquad\text{for every }e\in E.
\]

Finally, put
\[
z=\frac{x_++2y_+}{3}.
\]
Since $x_+,y_+\in E$ are isotropic and orthogonal,
\[
z\in\operatorname{conv}(E)\cap Q.
\]
On the other hand, its coordinates in the chosen basis give
\[
A_+(z)=\frac13,
\qquad
C_+(z)=\frac23.
\]
Therefore
\[
F(z)
=
\left(\frac{20}{3}-6\right)
\left(1-\frac43\right)
=
-\frac29<0.
\]
Consequently $z\notin E$, proving
\[
E\subsetneq\operatorname{conv}(E)\cap Q.
\]
\end{proof}

\begin{remark}
This example uses a generalized geometric realization in
which the entries associated with infinite Coxeter exponents
are equal to $-2$.
It does not establish the same conclusion for the canonical
Tits realization, where those entries are fixed at $-1$.
Moreover, the example concerns the full set of limit roots;
it does not by itself produce a strict dominance chain with
more than one accumulation point.
\end{remark}
\begin{proposition}\label{prop:nonexposed-extreme-limit-roots}
There exists an irreducible based root system $\Phi$ of rank eight
and signature $(5,3)$ such that
\[
E_{\mathrm{exp}}(\Phi)
\subsetneq E_{\mathrm{ext}}(\Phi).
\]
More precisely, its limit-root convex hull has two distinct
extreme points $x,y$ such that every supporting hyperplane at
either point contains both points.
\end{proposition}

\begin{proof}
Let $V=\bigoplus_{i=1}^{8}\mathbb R\alpha_i$,
and define $B$ by the Gram matrix
\[
G=
\begin{pmatrix}
1&-2&0&0&0&0&0&0\\
-2&1&-2&0&0&0&0&0\\
0&-2&1&-2&0&0&0&0\\
0&0&-2&1&-2&0&0&0\\
0&0&0&-2&1&-2&0&0\\
0&0&0&0&-2&1&-2&0\\
0&0&0&0&0&-2&1&-2\\
0&0&0&0&0&0&-2&1
\end{pmatrix}.
\]
Let
\[
s_i(v)=v-2B(v,\alpha_i)\alpha_i,
\qquad
W=\langle s_1,\ldots,s_8\rangle,
\qquad
\Phi=W\{\alpha_1,\ldots,\alpha_8\}.
\]
These data define a based root system, and its Coxeter graph
is connected.
The eigenvalues of $G$ are
\[
1-4\cos\frac{k\pi}{9},
\qquad 1\le k\le8.
\]
Exactly three are negative, so the signature is $(5,3)$.

Use the normalization
\[
h\left(\sum_i c_i\alpha_i\right)=\sum_i c_i,
\qquad
\widehat v=\frac{v}{h(v)}.
\]
Write
\[
E=E(\Phi),
\qquad
C=\operatorname{conv}(E),
\qquad
K=\mathbb R_{\ge0}C.
\]
The cone $K$ is closed and $W$-invariant, is contained in the
nonnegative simple-root cone, and satisfies
\[
B(v,w)\le0
\qquad(v,w\in K).
\]
These are standard properties of the closed imaginary cone;
see \cite[Section~2]{DHR}.

\medskip
\noindent\textit{The three dihedral blocks.}
Put
\[
(a_1,b_1)=(1,2),
\qquad
(a_2,b_2)=(4,5),
\qquad
(a_3,b_3)=(7,8),
\]
and set
\[
\tau=2-\sqrt3,
\qquad
\lambda=\tau^{-2}=7+4\sqrt3.
\]
For $1\le i\le3$, define
\[
\xi_i=
\frac{\alpha_{a_i}+\tau\alpha_{b_i}}{1+\tau},
\qquad
\eta_i=
\frac{\tau\alpha_{a_i}+\alpha_{b_i}}{1+\tau}.
\]
These are the two limit roots of the corresponding non-affine
dihedral subsystem. In particular,
\[
\xi_i,\eta_i\in E,
\qquad
B(\xi_i,\eta_i)=-1.
\]
The three two-dimensional blocks are pairwise orthogonal.

We will use that $\xi_1$ and $\xi_3$ generate extreme rays of $K$.
Indeed, for $i\in\{1,3\}$, the functional
\[
v\longmapsto
B(\xi_i,v)-\sum_{j\notin\{a_i,b_i\}}c_j(v)
\]
is nonpositive on $C$.
Equality forces $v$ to lie in the corresponding dihedral block
and to be orthogonal to $\xi_i$.
Since that block is nondegenerate of signature $(1,1)$,
its intersection with $\xi_i^\perp$ is $\mathbb R\xi_i$.
The normalization $h(v)=1$ then gives $v=\xi_i$.
Thus $\xi_1$ and $\xi_3$ are exposed points of $C$.

Define
\[
a=s_1s_2,
\qquad
b=s_4s_5,
\qquad
c=s_7s_8,
\qquad
T=s_3s_6.
\]
The elements $a,b,c$ commute. On their respective blocks,
they have eigenvalues $\lambda,\lambda^{-1}$ with
eigenvectors $\xi_i,\eta_i$.

\medskip
\noindent\textit{A family approaching a rank-two projection.}
Define
\[
P_i^+(v)=-B(v,\eta_i)\xi_i,
\qquad
P_i^-(v)=-B(v,\xi_i)\eta_i,
\]
and
\[
P_U=I-\sum_{i=1}^{3}(P_i^++P_i^-).
\]
Here $P_U$ is the projection onto the orthogonal complement of
the sum of the three dihedral blocks.

Put $\gamma=7/3$, and define
\[
\begin{aligned}
A_0&=\gamma^{-1}(P_1^++P_3^+)T,\\
A_1&=\gamma^{-1}P_2^+T,\\
A_2&=\gamma^{-1}P_UT,\\
A_3&=\gamma^{-1}P_2^-T,\\
A_4&=\gamma^{-1}(P_1^-+P_3^-)T.
\end{aligned}
\]
For every integer $d\ge1$, set
\[
g(d)=a^{2d}b^dc^{2d}T,
\qquad
F(d)=\frac{g(d)}{\gamma\lambda^{2d}}.
\]
The spectral decomposition of the three commuting blocks gives
the exact identity
\begin{equation}\label{eq:nonexposed-factor-expansion}
F(d)=
A_0+\lambda^{-d}A_1+\lambda^{-2d}A_2
+\lambda^{-3d}A_3+\lambda^{-4d}A_4.
\end{equation}

Direct calculation gives
\begin{equation}\label{eq:nonexposed-projection}
A_0^2=A_0,
\qquad
\operatorname{im}A_0=\operatorname{span}\{\xi_1,\xi_3\},
\qquad
A_0\xi_1=\xi_1,
\qquad
A_0\xi_3=\xi_3.
\end{equation}
Moreover, with
\[
u=\frac{8+4\sqrt3}{7},
\qquad
v=\frac{8-4\sqrt3}{7},
\]
we have $u,v>0$ and
\begin{equation}\label{eq:nonexposed-transitions}
A_0\xi_2=u\xi_1+v\xi_3,
\qquad
A_1\xi_1=v\xi_2,
\qquad
A_1\xi_3=u\xi_2.
\end{equation}
For example, the matrix of $(P_1^++P_2^++P_3^+)T$
on the ordered vectors $(\xi_1,\xi_2,\xi_3)$ is
\[
\begin{pmatrix}
\frac73&\frac{8+4\sqrt3}{3}&0\\
\frac{8-4\sqrt3}{3}&\frac{11}{3}&\frac{8+4\sqrt3}{3}\\
0&\frac{8-4\sqrt3}{3}&\frac73
\end{pmatrix},
\]
from which these identities follow.

\medskip
\noindent\textit{Convergence of the products.}
Set
\[
d_n=n+10,
\qquad
\varepsilon_n=\lambda^{-d_n},
\qquad
F_n=F(d_n),
\qquad
L_n=F_1\cdots F_n.
\]
Each $L_n$ is a positive scalar multiple of a group element,
and hence
\[
L_nK=K,
\qquad
L_n^{-1}K=K.
\]

For completeness, we verify convergence in operator norm.
Equip $V$ with the norm
\[
\|z\|_*=
\|A_0z\|_\infty+\|(I-A_0)z\|_\infty,
\]
where $\|\cdot\|_\infty$ is the maximum norm in the simple-root
coordinates. Since $A_0$ is a projection,
\[
\|A_0\|_*\le1,
\qquad
\|I-A_0\|_*\le1.
\]
Write
\[
F_n=A_0+R_n,
\qquad
\delta_n=\|R_n\|_*.
\]

The explicit matrices above give the elementary bounds
\[
\|A_0\|_\infty<6,
\qquad
\|A_j\|_\infty<8
\quad(1\le j\le4).
\]
Consequently,
\[
\|z\|_\infty\le\|z\|_*\le13\|z\|_\infty,
\]
and \eqref{eq:nonexposed-factor-expansion} implies, with a
deliberately coarse bound,
\[
\delta_n\le5000\varepsilon_n.
\]
In particular,
\begin{equation}\label{eq:nonexposed-small-perturbation}
S:=\sum_{n\ge1}\delta_n
\le
\frac{5000\lambda^{-11}}{1-\lambda^{-1}}
<10^{-8}.
\end{equation}

We have the uniform bound
\[
\|L_n\|_*
\le
\prod_{j=1}^{n}(1+\delta_j)
\le e^S.
\]
Also,
\[
L_n(I-A_0)=L_{n-1}R_n(I-A_0),
\]
so
\[
\|L_n(I-A_0)\|_*\le e^S\delta_n.
\]
It follows that
\[
\begin{aligned}
L_{n+1}-L_n
&=L_nR_{n+1}-L_n(I-A_0),\\
\|L_{n+1}-L_n\|_*
&\le e^S(\delta_{n+1}+\delta_n).
\end{aligned}
\]
Thus $(L_n)$ converges to an operator $L$.
The preceding estimate also yields
\[
L(I-A_0)=0,
\qquad\text{hence}\qquad
L=LA_0.
\]

Comparing the products $F_1\cdots F_n$ and $A_0^n=A_0$
by telescoping gives
\[
\|L-A_0\|_*\le e^S-1<1.
\]
Therefore $L$ is injective on $\operatorname{im}A_0$.
Together with $L=LA_0$, this proves
\[
\operatorname{rank}L=2.
\]

Applying the same argument to each tail gives operators
\[
Q_n=\lim_{m\to\infty}F_{n+1}\cdots F_m
\]
such that
\begin{equation}\label{eq:nonexposed-tail}
L=L_nQ_n,
\qquad
Q_n\longrightarrow A_0.
\end{equation}

\medskip
\noindent\textit{The two limiting extreme rays.}
Put
\[
X=L\xi_1,
\qquad
Y=L\xi_3,
\qquad
x=\widehat X,
\qquad
y=\widehat Y.
\]
The vectors $X,Y$ belong to $K$, are nonzero, and are linearly
independent. In particular, $h(X),h(Y)>0$.
Since normalized group actions preserve $E$,
\[
\widehat{L_n\xi_1},\widehat{L_n\xi_3}\in E.
\]
Passing to the limit gives $x,y\in E$, with $x\ne y$.

We show that $X$ generates an extreme ray of $K$.
Suppose
\[
X=A+B,
\qquad A,B\in K.
\]
For every $n$, put
\[
A_n=L_n^{-1}A,
\qquad
B_n=L_n^{-1}B.
\]
These vectors belong to $K$, and
\[
A_n+B_n=L_n^{-1}X=Q_n\xi_1\longrightarrow\xi_1.
\]
Because $K$ is contained in the nonnegative coordinate cone,
both sequences $(A_n)$ and $(B_n)$ are bounded.
After passing to a subsequence, write
\[
A_n\longrightarrow A_\infty,
\qquad
B_n\longrightarrow B_\infty.
\]
Then
\[
A_\infty,B_\infty\in K,
\qquad
A_\infty+B_\infty=\xi_1.
\]
Since $\xi_1$ generates an extreme ray,
$A_\infty=t\xi_1$ for some $t\ge0$.
Therefore
\[
A=L_nA_n\longrightarrow L(t\xi_1)=tX.
\]
Thus every decomposition of $X$ in $K$ is along its ray.
The same argument applies to $Y$, and hence
\begin{equation}\label{eq:nonexposed-extremality}
x,y\in E_{\mathrm{ext}}(\Phi).
\end{equation}

\medskip
\noindent\textit{The supporting hyperplanes cannot separate the endpoints.}
Set
\[
X_n=L_n\xi_1,
\qquad
Y_n=L_n\xi_3,
\qquad
Z=L\xi_2.
\]
Since $L=LA_0$, equation
\eqref{eq:nonexposed-transitions} gives
\begin{equation}\label{eq:nonexposed-mixed-direction}
Z=uX+vY.
\end{equation}

Using \eqref{eq:nonexposed-factor-expansion} and
\eqref{eq:nonexposed-transitions}, together with the uniform
boundedness of $L_n$, we obtain
\[
X_{n+1}-X_n
=
v\varepsilon_{n+1}L_n\xi_2
+O(\varepsilon_{n+1}^2).
\]
Put
\[
t_n=\sum_{k>n}\varepsilon_k.
\]
Since $L_n\xi_2\to Z$ and
$\sum_{k>n}\varepsilon_k^2=o(t_n)$, summing the preceding
identity gives
\[
X-X_n=vt_nZ+o(t_n).
\]
By \eqref{eq:nonexposed-mixed-direction},
\[
X_n=(1-uvt_n)X-v^2t_nY+o(t_n).
\]
After normalization, this becomes
\begin{equation}\label{eq:nonexposed-tangent-x}
\widehat{X_n}
=
x+
v^2\frac{h(Y)}{h(X)}t_n(x-y)
+o(t_n).
\end{equation}

Let $\ell$ be any linear functional attaining its maximum
over $C$ at $x$.
As $\widehat{X_n}\in E\subseteq C$,
\[
\ell(\widehat{X_n})-\ell(x)\le0.
\]
Divide \eqref{eq:nonexposed-tangent-x}, after applying $\ell$,
by $t_n$ and let $n\to\infty$.
Because $v^2h(Y)/h(X)>0$, this gives
\[
\ell(x)\le\ell(y).
\]
The maximality at $x$ gives the reverse inequality, so
\[
\ell(x)=\ell(y).
\]
Since $x\ne y$, the point $x$ is not exposed.

Likewise,
\[
Y_{n+1}-Y_n
=
u\varepsilon_{n+1}L_n\xi_2
+O(\varepsilon_{n+1}^2),
\]
and therefore
\[
\widehat{Y_n}
=
y+
u^2\frac{h(X)}{h(Y)}t_n(y-x)
+o(t_n).
\]
The same argument shows that every linear functional attaining
its maximum at $y$ also attains that maximum at $x$.
Thus $y$ is not exposed either.

Combining this with \eqref{eq:nonexposed-extremality} yields
\[
x,y\in
E_{\mathrm{ext}}(\Phi)\setminus E_{\mathrm{exp}}(\Phi),
\]
as required.
\end{proof}

\begin{remark}
The counterexample points are specified by the group elements
\[
g_n=
(s_1s_2)^{2(n+10)}
(s_4s_5)^{n+10}
(s_7s_8)^{2(n+10)}
s_3s_6:
\]
namely,
\[
x=\lim_{n\to\infty}
\widehat{g_1\cdots g_n\xi_1},
\qquad
y=\lim_{n\to\infty}
\widehat{g_1\cdots g_n\xi_3}.
\]
The summable perturbations preserve two extreme directions,
while their first-order terms force every supporting hyperplane
at either limiting point to contain the other point.
\end{remark}

\begin{remark}
All infinite edges in this example have Gram parameter $-2$.
Replacing these entries by $-1$ changes the spectral
decomposition used in the proof: the corresponding rank-two
blocks become affine.
Consequently, the proposition does not answer whether
\[
E_{\mathrm{exp}}(\Phi)=E_{\mathrm{ext}}(\Phi)
\]
holds for every standard Tits realization.
\end{remark}
\begin{theorem}\label{thm:orthogonal-parabolic-join}
Assume that $(\Phi,\Pi)$ is irreducible. Let $I,J\subseteq S$
determine irreducible infinite standard parabolic subsystems
satisfying
\[
B(V_I,V_J)=0.
\]
If $\Phi_I$ and $\Phi_J$ are not both non-affine root systems of
rank two, then
\[
E_I*E_J\subseteq E.
\]

For subsets $A,B\subseteq\{v:h(v)=1\}$, define their join by
\[
A*B
=
\{(1-t)x+ty:x\in A,\ y\in B,\ 0\le t\le1\}.
\]

The excluded case cannot be omitted in general: there exists an
irreducible rank-five based root system with two orthogonal
non-affine rank-two standard parabolic subsystems for which
\[
E_I*E_J\not\subseteq E.
\]
\end{theorem}

\begin{remark}\label{rem:meaning-of-exception}
The exception concerns two simultaneously non-affine rank-two
root systems. It does not concern the property of being parabolic:
both subsystems in the theorem are standard parabolic subsystems.

Moreover, the theorem does not assert that the join inclusion
always fails in the excluded case. It asserts that the inclusion
holds outside that case and can fail within it.
\end{remark}

We use two results of Dyer--Hohlweg--Ripoll. First, if a based root
system is irreducible, every orbit in its limit-root set is dense
\cite[Theorem~3.1]{DHR}. Second, an irreducible infinite non-affine
based root system contains a based root subsystem whose simple
roots are positive roots, span the same ambient space, and have
pairwise inner products strictly less than $-1$
\cite[Proposition~6.4]{DHR}.

\begin{lemma}\label{lem:connecting-root}
Under the hypotheses of Theorem~\ref{thm:orthogonal-parabolic-join},
there exists $\rho\in\Phi^+$ such that
\[
\operatorname{supp}(\rho)\subseteq S\setminus(I\cup J),
\]
\[
B(\rho,\alpha)\le0
\qquad
(\alpha\in\Phi_I^+\cup\Phi_J^+),
\]
and neither $B(\rho,-)|_{V_I}$ nor $B(\rho,-)|_{V_J}$ is zero.
Moreover,
\[
w\rho\in\Phi^+
\qquad (w\in W_IW_J).
\]
\end{lemma}

\begin{proof}
Choose a shortest path in the Coxeter graph joining a vertex of
$I$ to a vertex of $J$. Since $B(V_I,V_J)=0$, this path has at least
one internal vertex. By minimality, all its internal vertices lie
outside $I\cup J$.

Starting with the simple root at the first internal vertex,
successively reflect at the remaining internal vertices along
the path. At each step the new simple root has strictly negative
inner product with the root already constructed, so the resulting
root is positive and acquires a strictly positive coefficient at
the new vertex. We obtain a positive root $\rho$ supported
precisely on the internal vertices.

All off-diagonal entries of the Gram matrix are nonpositive.
Consequently,
\[
B(\rho,\alpha)\le0
\qquad
(\alpha\in\Phi_I^+\cup\Phi_J^+).
\]
The first and last edges of the chosen path show that the
restrictions of $B(\rho,-)$ to $V_I$ and $V_J$ are nonzero.

Every element of $W_IW_J$ preserves the coordinates outside
$I\cup J$. At least one such coordinate of $\rho$ is strictly
positive. Hence $w\rho$ cannot be a negative root. The sign
property of roots therefore gives $w\rho\in\Phi^+$.
\end{proof}

\begin{lemma}\label{lem:one-segment-suffices}
Suppose that $x\in E_I$ and $y\in E_J$ satisfy
\[
[x,y]\subseteq E.
\]
Then
\[
E_I*E_J\subseteq E.
\]
\end{lemma}

\begin{proof}
The group $W_I$ fixes $V_J$ pointwise, and $W_J$ fixes $V_I$
pointwise. Thus, for $u\in W_I$ and $v\in W_J$,
\[
(uv)\cdot[x,y]=[u\cdot x,v\cdot y]\subseteq E.
\]
Here the normalized linear action maps the segment onto the
segment between the normalized images of its endpoints.

By minimality,
\[
\overline{W_I\cdot x}=E_I,
\qquad
\overline{W_J\cdot y}=E_J.
\]
Given $x'\in E_I$ and $y'\in E_J$, choose sequences
$u_n\in W_I$ and $v_n\in W_J$ such that
\[
u_n\cdot x\longrightarrow x',
\qquad
v_n\cdot y\longrightarrow y'.
\]
For every $t\in[0,1]$,
\[
(1-t)(u_n\cdot x)+t(v_n\cdot y)\in E.
\]
The closedness of $E$ gives
\[
(1-t)x'+ty'\in E.
\]
\end{proof}

\begin{proof}[Proof of the inclusion in
Theorem~\ref{thm:orthogonal-parabolic-join}]
Fix the positive root $\rho$ supplied by
Lemma~\ref{lem:connecting-root}. By
Lemma~\ref{lem:one-segment-suffices}, it is enough to construct one
segment $[x,y]\subseteq E$ with $x\in E_I$ and $y\in E_J$.

\medskip
\noindent
\textit{Case 1: at least one subsystem is affine.}

Suppose first that $\Phi_I$ is affine. Let
\[
E_I=\{\delta\},
\qquad h(\delta)=1.
\]
The affine root-string structure provides a root $a$ and a
constant $\ell>0$ such that
\[
a_n=a+n\ell\delta\in\Phi_I^+
\qquad(n\ge0).
\]
Since $\delta$ has strictly positive coordinates on $I$,
Lemma~\ref{lem:connecting-root} implies
\[
B(\rho,\delta)<0.
\]
Consequently,
\[
\begin{aligned}
s_{a_n}\rho
&=\rho-2B(\rho,a_n)a_n\\
&=Cn^2\delta+O(n),
\end{aligned}
\]
where
\[
C=-2\ell^2B(\rho,\delta)>0.
\]

Assume that $\Phi_J$ is non-affine. By
\cite[Proposition~6.4]{DHR}, we can choose positive roots
$c,d\in\Phi_J^+$ such that
\[
B(c,d)<-1,
\qquad B(\rho,c)<0.
\]
Let $t=s_cs_d$, and let $y\in E_J$ be its normalized attracting
eigenvector, with eigenvalue $\mu>1$. Since
\[
B(\rho,c)<0,\qquad B(\rho,d)\le0,
\]
the attracting coefficient of $\rho$ is strictly positive.
Thus, for some $D>0$,
\[
t^m\rho=D\mu^m y+O(1).
\]
The element $t$ fixes $V_I$ pointwise, so
\[
s_{a_n}t^m\rho
=
Cn^2\delta+D\mu^m y+O(n+1).
\]
All these vectors are positive roots by
Lemma~\ref{lem:connecting-root}.

Given $r>0$, choose positive integers $n_m$ satisfying
\[
n_m\sim\sqrt{\frac{rD}{C}\mu^m}.
\]
Then
\[
\widehat{s_{a_{n_m}}t^m\rho}
\longrightarrow
\frac{r\delta+y}{r+1}.
\]
Hence $[\delta,y]\subseteq E$.

If $\Phi_J$ is also affine, write
\[
E_J=\{\varepsilon\},
\qquad h(\varepsilon)=1,
\]
and choose an affine root string
\[
b_m=b+m\ell'\varepsilon\in\Phi_J^+.
\]
Orthogonality gives
\[
s_{a_n}s_{b_m}\rho
=
Cn^2\delta+Dm^2\varepsilon+O(n+m+1)
\]
for constants $C,D>0$.
For any $r>0$, choosing $n,m\to\infty$ with
\[
\frac{Cn^2}{Dm^2}\longrightarrow r
\]
yields
\[
\widehat{s_{a_n}s_{b_m}\rho}
\longrightarrow
\frac{r\delta+\varepsilon}{r+1}.
\]
Thus $[\delta,\varepsilon]\subseteq E$, completing this case.

\medskip
\noindent
\textit{Case 2: both subsystems are non-affine.}

By hypothesis, one subsystem has rank at least three.
Interchanging $I$ and $J$ if necessary, assume
\[
\dim V_I\ge3.
\]
Using \cite[Proposition~6.4]{DHR}, choose linearly independent
positive roots $a,b,u\in\Phi_I^+$ that belong to the simple system
of a based root subsystem and satisfy
\[
B(a,b),B(a,u),B(b,u)<-1,
\qquad B(\rho,a)<0.
\]
Likewise choose $c,d\in\Phi_J^+$ with
\[
B(c,d)<-1,\qquad B(\rho,c)<0.
\]

Set
\[
g=s_as_b,\qquad t=s_cs_d.
\]
Choose normalized eigenvectors
\[
gx=\lambda x,\qquad gx_-=\lambda^{-1}x_-,
\]
\[
ty=\mu y,\qquad ty_-=\mu^{-1}y_-,
\]
where
\[
\lambda,\mu>1,\qquad
x,x_-\in E_I,\qquad y,y_-\in E_J.
\]
The planes
\[
L=\operatorname{span}\{a,b\},
\qquad
M=\operatorname{span}\{c,d\}
\]
are nondegenerate and mutually orthogonal. Every $v\in V$
therefore has a unique decomposition
\[
v=A(v)x+A_-(v)x_-+D(v)y+D_-(v)y_-+v_0,
\]
where
\[
v_0\in(L\oplus M)^\perp.
\]
The connecting-root inequalities imply
\[
A(\rho),A_-(\rho),D(\rho),D_-(\rho)>0.
\]

Suppose first that
\[
\frac{\log\lambda}{\log\mu}\notin\mathbb Q.
\]
For every $r>0$, irrational rotation on the circle gives sequences
$n_k,m_k\to\infty$ such that
\[
\frac{A(\rho)\lambda^{n_k}}
     {D(\rho)\mu^{m_k}}
\longrightarrow r.
\]
Since
\[
g^{n_k}t^{m_k}\rho
=
A(\rho)\lambda^{n_k}x
+D(\rho)\mu^{m_k}y+O(1),
\]
we obtain
\[
\widehat{g^{n_k}t^{m_k}\rho}
\longrightarrow
\frac{rx+y}{r+1}.
\]
Thus $[x,y]\subseteq E$.

It remains to consider
\[
\frac{\log\lambda}{\log\mu}\in\mathbb Q.
\]
Replacing $g$ and $t$ by suitable positive powers, we may assume
that their attracting eigenvalues coincide:
\[
\lambda=\mu=\Lambda>1.
\]

Define
\[
\mathcal R
=
\left\{
r>0:\frac{rx+y}{r+1}\in E
\right\}.
\]
This is a closed subset of $(0,\infty)$. Since $g$ fixes $y$ and
sends $x$ to $\Lambda x$, the $W$-invariance of $E$ gives
\[
\Lambda^{\mathbb Z}\mathcal R=\mathcal R.
\]

Put
\[
\kappa=-B(x,x_-)>0.
\]
The inequalities $B(u,a),B(u,b)<-1$ imply
\[
A(u)>0,\qquad A_-(u)>0.
\]
Set
\[
d_0=\kappa A(u)A_-(u)>0,
\qquad
u_m=g^m u.
\]
Because $a,b,u$ belong to a simple system of a based root
subsystem, $u_m\in\Phi_I^+$ for $m\ge0$.

Writing
\[
u=A(u)x+A_-(u)x_-+u_0,
\qquad u_0\perp L,
\]
and using $B(u,u)=1$, we obtain
\[
B(u,u_m)
=
1+2d_0-d_0(\Lambda^m+\Lambda^{-m}).
\]
For all sufficiently large $m$, therefore,
\[
q_m:=-B(u,u_m)>1.
\]
Let
\[
H_m=s_{u_m}s_u.
\]
Its attracting eigenvalue is
\[
\eta_m=
\left(q_m+\sqrt{q_m^2-1}\right)^2.
\]

For sufficiently large $m$,
\[
B(\rho,u)\le0,\qquad B(\rho,u_m)<0.
\]
Both isotropic eigenvectors of $H_m$ in
$\operatorname{span}\{u,u_m\}$ are positive linear combinations
of $u$ and $u_m$. It follows that the attracting coefficient of
$\rho$ is strictly positive.

Moreover,
\[
A(u)>0,\qquad A(u_m)=\Lambda^mA(u)>0.
\]
Applying the linear functional $A$ to the eigenspace decomposition
of $H_m^k\rho$ consequently gives
\[
A(H_m^k\rho)=C_m\eta_m^k+O(1)
\qquad(k\to\infty)
\]
for some $C_m>0$.

Since $H_m\in W_I$,
\[
D(H_m^k\rho)=D(\rho)=:D_0>0.
\]
Fixing a sufficiently large $k$ and letting $N\to\infty$, we have
\[
\widehat{g^Nt^NH_m^k\rho}
\longrightarrow
\frac{A(H_m^k\rho)x+D_0y}
     {A(H_m^k\rho)+D_0}.
\]
Hence
\[
\frac{A(H_m^k\rho)}{D_0}\in\mathcal R.
\]

Consider the closed subset
\[
\mathcal T
=
\{\log_\Lambda r+\mathbb Z:r\in\mathcal R\}
\subseteq\mathbb R/\mathbb Z.
\]
Its closedness follows from the closedness of $\mathcal R$ and
the invariance $\Lambda^{\mathbb Z}\mathcal R=\mathcal R$.
For every sufficiently large $m$, the set $\mathcal T$ contains
all accumulation points of
\[
k\log_\Lambda\eta_m
+\log_\Lambda(C_m/D_0)+o(1)
\pmod{\mathbb Z}.
\]

If some $\log_\Lambda\eta_m$ is irrational, these accumulation
points fill the circle, and hence
\[
\mathcal T=\mathbb R/\mathbb Z.
\]
If all $\log_\Lambda\eta_m$ are rational but their reduced
denominators are unbounded, then $\mathcal T$ contains translates
of equally spaced grids of arbitrarily large cardinality.
Closedness again gives
\[
\mathcal T=\mathbb R/\mathbb Z.
\]

Suppose, finally, that all these numbers are rational with
uniformly bounded reduced denominators. Then there exists
$M\in\mathbb N$ such that
\[
\log_\Lambda\eta_m-2m\in\frac1M\mathbb Z
\]
for every sufficiently large $m$.

On the other hand,
\[
q_m=d_0\Lambda^m-(1+2d_0)+d_0\Lambda^{-m},
\]
so
\[
\frac{\eta_m}{\Lambda^{2m}}\longrightarrow4d_0^2.
\]
Furthermore,
\[
q_m<d_0\Lambda^m,
\qquad
\eta_m<4q_m^2,
\]
and therefore
\[
\frac{\eta_m}{\Lambda^{2m}}<4d_0^2.
\]
Thus the sequence
\[
\log_\Lambda\eta_m-2m
\]
lies in the discrete set $\frac1M\mathbb Z$, converges to
$\log_\Lambda(4d_0^2)$, and is strictly smaller than its limit.
This is impossible.

Consequently,
\[
\mathcal T=\mathbb R/\mathbb Z.
\]
Together with $\Lambda^{\mathbb Z}\mathcal R=\mathcal R$, this
implies
\[
\mathcal R=(0,\infty).
\]
Hence $[x,y]\subseteq E$. Applying
Lemma~\ref{lem:one-segment-suffices} completes the proof of the
join inclusion.
\end{proof}

\begin{proof}[Construction and verification of the counterexample]
Let $\Pi=\{\alpha_1,\ldots,\alpha_5\}$, and take the Gram matrix
\[
G=
\begin{pmatrix}
1&-2&0&0&-2\\
-2&1&0&0&-2\\
0&0&1&-2&-2\\
0&0&-2&1&-2\\
-2&-2&-2&-2&1
\end{pmatrix}.
\]
The Coxeter graph is connected, so the ambient root system is
irreducible. Set
\[
I=\{1,2\},\qquad J=\{3,4\}.
\]
Then $V_I\perp V_J$, and both standard parabolic subsystems are
non-affine of rank two.

Put
\[
\tau=2-\sqrt3,
\qquad
x=\frac{\alpha_1+\tau\alpha_2}{1+\tau},
\qquad
y=\frac{\alpha_3+\tau\alpha_4}{1+\tau}.
\]
The matrix of $s_1s_2$ on $V_I$ is
\[
\begin{pmatrix}
15&-4\\
4&-1
\end{pmatrix},
\]
whose attracting eigenvalue is $\tau^{-2}$ and whose attracting
eigenvector is $(1,\tau)$. Hence
\[
\widehat{(s_1s_2)^n\alpha_1}\longrightarrow x.
\]
Similarly,
\[
\widehat{(s_3s_4)^n\alpha_3}\longrightarrow y.
\]
Thus
\[
x\in E_I,\qquad y\in E_J.
\]

We shall prove that
\[
z=\frac{x+2y}{3}\notin E.
\]

Define
\[
p=\frac{1+\sqrt3}{2},
\qquad
q=\frac{\sqrt3-1}{2}.
\]
For
\[
v=c_1\alpha_1+c_2\alpha_2+c_3\alpha_3+c_4\alpha_4+d\alpha_5,
\]
set
\[
\begin{aligned}
A_+(v)&=pc_1-qc_2+2d,
&
A_-(v)&=-qc_1+pc_2+2d,\\
C_+(v)&=pc_3-qc_4+2d,
&
C_-(v)&=-qc_3+pc_4+2d.
\end{aligned}
\]
Consider the homogeneous quadratic polynomial
\[
H(v)
=
\bigl(20A_+(v)-9C_+(v)\bigr)
\bigl(3A_+(v)-2C_+(v)\bigr).
\]
We claim that
\[
H(\beta)\ge0
\qquad(\beta\in\Phi^+).
\]

All roots have integral coordinates in the basis $\Pi$.
If $\beta=\sum_i b_i\alpha_i$ is a positive root that is not
simple, then
\[
1=B(\beta,\beta)=\sum_i b_iB(\beta,\alpha_i)
\]
provides an index $i$ with
\[
b_i>0,\qquad B(\beta,\alpha_i)>0.
\]
The root $s_i\beta$ remains positive: another positive coordinate
of $\beta$ is unchanged, so $s_i\beta$ cannot be negative.
Moreover,
\[
s_i\beta=\beta-2B(\beta,\alpha_i)\alpha_i
\]
has strictly smaller height and differs from $\beta$ only in its
$i$-th coordinate.

Iterating this descent reaches a simple root, since the heights
are positive integers. Reversing the descent shows that every
positive root admits a construction from a simple root along
which all coordinates are nondecreasing.

Fix such a construction of $\beta$, and suppose first that its
fifth coordinate $d$ is positive. If the construction uses $s_5$,
let $\rho$ be the root immediately after the last application of
$s_5$. Write
\[
\rho=a_1\alpha_1+\cdots+a_4\alpha_4+d\alpha_5,
\qquad
S=a_1+\cdots+a_4.
\]
If $d_{\mathrm{old}}$ is the fifth coordinate immediately before
that reflection, then
\[
d=4S-d_{\mathrm{old}}.
\]
Coordinate monotonicity gives $d\ge d_{\mathrm{old}}$, and hence
\[
d\ge2S.
\]
If the construction never uses $s_5$, its initial simple root
must be $\alpha_5$; in this case take $\rho=\alpha_5$, and the same
inequality holds.

In either case,
\[
\beta=w\rho
\qquad\text{for some }w\in W_IW_J,
\]
and
\[
a_i\ge0,\qquad \sum_{i=1}^4a_i\le\frac d2.
\]
Consequently,
\[
Ld\le A_\pm(\rho),C_\pm(\rho)\le Ud,
\]
where
\[
L=2-\frac q2=\frac{9-\sqrt3}{4},
\qquad
U=2+\frac p2=\frac{9+\sqrt3}{4}.
\]
Set
\[
R=\frac UL=\frac{9+\sqrt3}{9-\sqrt3}<\frac32.
\]

A direct calculation gives
\[
\begin{aligned}
(A_+,A_-)(s_1v)
&=(\tau^{-1}A_-(v),\,\tau A_+(v)),\\
(A_+,A_-)(s_2v)
&=(\tau A_-(v),\,\tau^{-1}A_+(v)).
\end{aligned}
\]
Both $s_1$ and $s_2$ fix $C_+$ and $C_-$. Likewise,
\[
\begin{aligned}
(C_+,C_-)(s_3v)
&=(\tau^{-1}C_-(v),\,\tau C_+(v)),\\
(C_+,C_-)(s_4v)
&=(\tau C_-(v),\,\tau^{-1}C_+(v)),
\end{aligned}
\]
and $s_3,s_4$ fix $A_+,A_-$.

It follows that
\[
A_+(\beta),C_+(\beta)>0,
\qquad
\frac{A_+(\beta)}{C_+(\beta)}
\in
\bigcup_{m\in\mathbb Z}\tau^m[R^{-1},R].
\]
Since $\tau<3/10$ and $R<3/2$, we have
\[
\tau R<\frac9{20},
\qquad
R^{-1}>\frac23.
\]
For $m\ge1$,
\[
\tau^mR\le\tau R<\frac9{20},
\]
whereas for $m\le0$,
\[
\tau^mR^{-1}\ge R^{-1}>\frac23.
\]
Thus the displayed union avoids the interval
$(9/20,2/3)$. Therefore
\[
\left(
20\frac{A_+(\beta)}{C_+(\beta)}-9
\right)
\left(
3\frac{A_+(\beta)}{C_+(\beta)}-2
\right)
\ge0,
\]
which proves $H(\beta)\ge0$ when $d>0$.

Now suppose that the fifth coordinate of $\beta$ is zero.
Its coordinatewise nondecreasing construction cannot use $s_5$.
Since $W_I$ and $W_J$ preserve their respective summands, the
root $\beta$ belongs to $\Phi_I$ or $\Phi_J$. In the first case
$C_+(\beta)=0$, and in the second case $A_+(\beta)=0$. Accordingly,
\[
H(\beta)=60A_+(\beta)^2
\qquad\text{or}\qquad
H(\beta)=18C_+(\beta)^2.
\]
This proves the claim for every positive root.

The homogeneity of $H$ implies
\[
H(\widehat\beta)=\frac{H(\beta)}{h(\beta)^2}\ge0.
\]
Taking accumulation points and using continuity, we obtain
\[
E\subseteq\{v\in V:H(v)\ge0\}.
\]

Finally, direct substitution gives
\[
A_+(x)=1,\qquad C_+(x)=0,
\qquad
A_+(y)=0,\qquad C_+(y)=1.
\]
Hence
\[
A_+(z)=\frac13,\qquad C_+(z)=\frac23,
\]
and
\[
H(z)
=
\left(\frac{20}{3}-6\right)
\left(1-\frac43\right)
=
-\frac29<0.
\]
Therefore $z\notin E$, although $z\in[x,y]$ with
$x\in E_I$ and $y\in E_J$. This proves
\[
E_I*E_J\not\subseteq E.
\]
\end{proof}

\bibliographystyle{unsrt}
\bibliography{reference}
\end{document}